% پیش‌نویس مستقل؛ تا زمان تأیید در فصل اصلی فراخوانی نشود.
% بخش ۳.۸ فعلی دست‌نخورده است. برچسب‌های آن برای ادغام نهایی حفظ شده‌اند.
\section{جدایی و تمایز میان‌طبقه‌ای}
\label{cmp:sec:ch3-center-separation}

در بخش~\ref{cmp:sec:ch3-quality-criteria} تراکم و گستردگی
بازنمایی‌های اعضای هر طبقه را بررسی کردیم. این معیارها
ساختار هندسی درون طبقه را توصیف می‌کنند، اما موقعیت آن
نسبت به سایر طبقه‌ها را مستقیماً نشان نمی‌دهند. حتی اگر
بازنمایی‌های هر طبقه در یک نقطه متمرکز شوند، ممکن است
مراکز طبقه‌های مختلف بر یکدیگر منطبق باشند.

اهمیت جدایی مراکز از قاعدهٔ تصمیم مدل ناشی می‌شود.
مراکز حسابی تعریف‌شده در رابطه~\eqref{cmp:eq:ch3-class-center}
صرفاً ابزار توصیف بازنمایی‌ها نیستند، بلکه همان توزیع‌های
مرجعی هستند که ناهمانندی بازنمایی با آن‌ها در تعیین حاشیهٔ
تعلق به طبقه به کار می‌رود
(رابطه‌های~\eqref{cmp:eq:ch3-pairwise-margin}
و~\eqref{cmp:eq:ch3-fixed-representation-margin}).
ازاین‌رو، افزون بر نزدیکی اعضای هر طبقه به مرکز خود،
تمایز آن مرکز از مراکز رقیب نیز اهمیت دارد. اگر دو
مرکز منطبق باشند، حاشیهٔ دوتایی متناظر برای هر بازنمایی
ثابت، هرگاه تعریف‌شده باشد، صفر خواهد بود.

شکل~\ref{cmp:fig:ch3-compactness-separation-four-cases}
چهار وضعیت را در یک مثال هندسی نشان می‌دهد. در سطر
نخست پراکندگی کم و در سطر دوم پراکندگی زیاد است.
در ستون چپ مراکز به هم نزدیک‌تر و در ستون راست
از هم دورترند. در هر سطر تنها جدایی مراکز و در
هر ستون تنها مقیاس پراکندگی تغییر می‌کند.

\begin{figure}[!htbp]
\centering
\begin{LTR}
\centering
\resizebox{0.92\linewidth}{!}{%
\begin{tikzpicture}[
 x=1cm,y=1cm,
 panel/.style={draw=black!25,rounded corners=2pt},
 ptA/.style={fill=blue!72!black,draw=white,line width=0.25pt},
 ptB/.style={fill=orange!85!black,draw=white,line width=0.25pt},
 ctrA/.style={fill=blue!72!black,draw=black,line width=0.45pt},
 ctrB/.style={fill=orange!85!black,draw=black,line width=0.45pt},
 titlebox/.style={font=\small,align=center}
]
% مختصات نسبی متقارن‌اند؛ مربع‌ها میانگین حسابی واقعی نقاط‌اند.
\foreach \px/\py/\hsep/\sc/\pid in
 {0/3.55/0.27/0.43/1,5.65/3.55/1.30/0.43/2,
  0/0/0.27/1.15/3,5.65/0/1.30/1.15/4}{
 \begin{scope}[shift={(\px,\py)}]
  \draw[panel] (-2.65,-1.32) rectangle (2.65,1.32);
  \draw[black!50,densely dashed,line width=0.45pt]
       (0,-1.14)--(0,1.14);
  \draw[black!55,<->,line width=0.55pt] (-\hsep,0)--(\hsep,0);
  \foreach \dx/\dy in {-0.85/0.24,0.85/-0.24,
    -0.58/-0.52,0.58/0.52,-0.28/0.73,0.28/-0.73,
    -0.11/-0.30,0.11/0.30}{
    \filldraw[ptA] ({-\hsep+\sc*\dx},{\sc*\dy}) circle (1.8pt);
    \filldraw[ptB] ({\hsep-\sc*\dx},{-\sc*\dy}) circle (1.8pt);
  }
  \filldraw[ctrA] ({-\hsep-0.065},-0.065)
       rectangle ({-\hsep+0.065},0.065);
  \filldraw[ctrB] ({\hsep-0.065},-0.065)
       rectangle ({\hsep+0.065},0.065);
  \ifnum\pid=1\relax
   \node[titlebox] at (0,-1.57) {\rl{الف) پراکندگی کم، جدایی کم}};
  \fi
  \ifnum\pid=2\relax
   \node[titlebox] at (0,-1.57) {\rl{ب) پراکندگی کم، جدایی زیاد}};
  \fi
  \ifnum\pid=3\relax
   \node[titlebox] at (0,-1.57) {\rl{ج) پراکندگی زیاد، جدایی کم}};
  \fi
  \ifnum\pid=4\relax
   \node[titlebox] at (0,-1.57) {\rl{د) پراکندگی زیاد، جدایی زیاد}};
  \fi
 \end{scope}
}
\end{tikzpicture}%
}
\end{LTR}
\caption{چهار وضعیت تراکم درون‌طبقه‌ای و جدایی میان‌طبقه‌ای.
دایره‌ها بازنمایی اعضا و مربع‌ها مراکز حسابی دو طبقه‌اند.
در هر سطر پراکندگی ثابت است و در هر ستون جدایی مراکز ثابت
است. خط‌چین مرز نزدیک‌ترین مرکز در این مثال اقلیدسی
دوبعدی است، نه مرز محاسبه‌شدهٔ KL یا JSD.}
\label{cmp:fig:ch3-compactness-separation-four-cases}
\end{figure}

در هر دو قاب سمت راست، باوجود تفاوت پراکندگی، همهٔ
نقاط در سمت مرکز طبقهٔ خود قرار دارند. در مقابل،
جدایی کم مراکز می‌تواند حتی با پراکندگی کم نیز موجب
خطای تصمیم شود. بنابراین کاهش پراکندگی تمام مسئلهٔ
هندسی طبقه‌بندی را حل نمی‌کند و فشرده‌سازی بیشتر
طبقه‌ای که اعضایش به‌خوبی از طبقهٔ رقیب متمایزند،
الزاماً دقت را افزایش نمی‌دهد.

اکنون نخست چهار معیار ساده برای ارزیابی مراکز معرفی
می‌کنیم؛ سپس از مقایسهٔ مؤلفه‌های منتخب به سنجش تمایز
کل توزیع‌های مرکزی می‌رسیم.

\subsection{معیارهای انطباق مراکز با الگوهای مرجع}
\label{cmp:subsec:ch3-center-separation-criteria}
\label{cmp:subsubsec:ch3-T}
\label{cmp:subsubsec:ch3-BT}
\label{cmp:subsubsec:ch3-TNL}
\label{cmp:subsubsec:ch3-BTNL}

مطابق رابطه~\eqref{cmp:eq:ch3-class-center-prob}،
$\mathcal V_{cd}$ احتمال مشاهدهٔ برچسب $d$ در
همسایگی مبدأیی از طبقهٔ $c$ است. با وزن‌های مثبت
$\omega_c$ با مجموع یک و مجموعهٔ جایگشت‌های
$\mathfrak S_K$، چهار معیار زیر را تعریف می‌کنیم:
\begin{align}
T(\mathcal V)&=\sum_c\omega_c\mathcal V_{cc}
=\operatorname{tr}\bigl(\operatorname{diag}(\Omega)\mathcal V\bigr),
\label{cmp:eq:ch3-T-on-V}\\
BT(\mathcal V)&=\max_{\sigma\in\mathfrak S_K}
\sum_c\omega_c\mathcal V_{c,\sigma(c)},
\label{cmp:eq:ch3-BT-on-V}\\
TNL(\mathcal V)&=-\sum_c\omega_c\log\mathcal V_{cc},
\label{cmp:eq:ch3-TNL-on-V}\\
BTNL(\mathcal V)&=\min_{\sigma\in\mathfrak S_K}
\left[-\sum_c\omega_c\log\mathcal V_{c,\sigma(c)}\right].
\label{cmp:eq:ch3-BTNL-on-V}
\end{align}
معیارهای $T$ و $BT$ میانگین حسابی احتمال‌های
منتخب را به‌ترتیب در الگوی هم‌برچسبی و بهترین
انتساب یک‌به‌یک می‌سنجند. معیارهای لگاریتمی،
با قرارداد $-\log0=+\infty$، هزینهٔ منفی
لگاریتم درست‌نمایی همین الگوهای فرضی‌اند؛
معادل واگرایی KL از بردارهای یک‌داغ مرجع به
مراکز، نه زیان پیش‌بینی نمونه‌ها.
کران‌ها، مقایسه‌ها و شرایط تساوی این چهار معیار
در بخش~\ref{app:ch3-assignment-comparisons}
از پیوست~\ref{app:model-geometry} آمده‌اند.

این معیارها انطباق با الگوهای مرجع را می‌سنجند،
نه تمایز کامل توزیع‌های مرکزی. اگرچه در یافتن
انتساب جایگشتی کل ماتریس بررسی می‌شود، امتیاز
نهایی هر چهار معیار فقط از یک مؤلفهٔ هر مرکز
به دست می‌آید. از آنجا که در قاعدهٔ تصمیم مدل
کل بازنمایی با کل توزیع مرکزی مقایسه می‌شود،
در ادامه معیاری را معرفی می‌کنیم که همهٔ
مؤلفه‌های مراکز را در نظر می‌گیرد.

\subsection{تمایز اطلاعاتی مراکز}
\label{cmp:subsubsec:ch3-I-criterion}

همان‌گونه که در بخش تراکم هزینهٔ تجمیع بازنمایی‌های
اعضای یک طبقه را در مرکز حسابی آن با KL اندازه
گرفتیم، اکنون هزینهٔ تجمیع خود مراکز طبقه‌ها در
یک توزیع مشترک را می‌سنجیم. میانگین وزن‌دار مراکز:
\begin{equation}\label{cmp:eq:ch3-global-center}
\overline{\mathcal V}
=\sum_{c=1}^{K}\omega_c\mathcal V_c
=\Omega\mathcal V.
\end{equation}
معیار تمایز اطلاعاتی عبارت است از:
\begin{equation}\label{cmp:eq:ch3-I-on-V}
\boxed{I(\mathcal V)=
\sum_{c=1}^{K}\omega_c
D_{\mathrm{KL}}\bigl(\mathcal V_c\|\overline{\mathcal V}\bigr).}
\end{equation}
این معیار همان واگرایی جنسن--شنون وزن‌دار چند
توزیع است. همهٔ مؤلفه‌های مراکز در آن نقش دارند و
نیازی به انتخاب برچسب مقصد اختصاصی نیست.
از هویت آنتروپی مخلوط، مطابق
رابطه~\eqref{cmp:eq:ch3-I-entropy} در پیوست هندسه:
\begin{equation*}
\boxed{I(\mathcal V)=H(\overline{\mathcal V})
-\sum_{c=1}^{K}\omega_c H(\mathcal V_c).}
\end{equation*}

\paragraph{تعبیر اطلاعاتی.}
طبقهٔ مبدأ $C_0$ را با احتمال $\omega_c$ انتخاب
می‌کنیم و سپس عضوی از آن طبقه را به‌صورت یکنواخت
برمی‌گزینیم. برچسب مشاهده‌شده در همسایگی آن را با
$C_\xi$ نشان می‌دهیم؛ مشاهده مطابق قرارداد گام
وزن‌دار رابطه~\eqref{cmp:eq:ch3-weighted-observation-position}
انجام می‌شود. ماتریس مراکز، ماتریس انتقال کانال
گسسته از طبقهٔ مبدأ به برچسب همسایگی است:
\begin{equation}\label{cmp:eq:ch3-I-joint-distribution}
\mathrm P(C_0=c,C_\xi=d)=\omega_c\mathcal V_{cd},
\qquad
\mathrm P(C_\xi=d)=\sum_{c=1}^{K}\omega_c\mathcal V_{cd}
=\overline{\mathcal V}(d).
\end{equation}
بنابراین، طبق تعریف اطلاعات متقابل \cite{CoverThomas2006}:
\begin{equation}\label{cmp:eq:ch3-I-mutual-info}
\mathcal I(C_0;C_\xi)=
\sum_{c,d}\omega_c\mathcal V_{cd}
\log\frac{\mathcal V_{cd}}{\overline{\mathcal V}(d)}
=I(\mathcal V).
\end{equation}
جمله‌های با احتمال صفر مطابق قرارداد KL تفسیر
می‌شوند. در نتیجه $I$ میزان کاهش عدم‌قطعیت دربارهٔ
طبقهٔ مبدأ با مشاهدهٔ یک برچسب همسایگی را می‌سنجد:
\begin{equation*}
\boxed{I(\mathcal V)=H(C_0)-H(C_0\mid C_\xi)
=H(\Omega)-H(C_0\mid C_\xi).}
\end{equation*}
و مطابق رابطه~\eqref{cmp:eq:ch3-I-bounds} داریم:
\begin{equation*}
\boxed{0\le I(\mathcal V)\le H(\Omega)\le\log K.}
\end{equation*}
با وزن‌های مثبت، $I=0$ دقیقاً هنگامی برقرار است
که تمام مراکز یکسان باشند و $I=H(\Omega)$ دقیقاً
در ماتریس‌های جایگشتی به دست می‌آید. اثبات این
نتایج در بخش~\ref{app:ch3-I-bounds-proof}
از پیوست~\ref{app:model-geometry} آمده است.
این حالت ایدئال لزوماً با سیاست انتخاب ویژگی روی
هر گراف دست‌یافتنی نیست.

ارتباط $I$ با $T$ و $BT$ در رابطه‌های
~\eqref{cmp:eq:ch3-I-vs-T} و~\eqref{cmp:eq:ch3-I-vs-BT}
از پیوست هندسه آمده است. با وزن‌های برابر،
افزایش $BT$ کران پایین $I$ را افزایش می‌دهد،
نه لزوماً مقدار خود $I$ را. همچنین نامساوی فانو
در رابطه~\eqref{cmp:eq:ch3-I-Fano} به تشخیص از
روی یک برچسب همسایگی مربوط است، نه پیش‌بینی از
روی کل بازنمایی.

\paragraph{ارتباط با تراکم مبتنی بر KL.}
طبق رابطه~\eqref{cmp:eq:ch3-R-class}، $R_c$
میانگین واگرایی بازنمایی‌های اعضای طبقهٔ $c$ از
مرکز آن است و $R_{\mathrm{KL}}$ صورت تجمیع‌شدهٔ
آن در رابطه~\eqref{cmp:eq:ch3-R-average} است.
در مقابل، $I$ میانگین واگرایی مراکز از میانگین
مشترک آن‌هاست. هویت تجزیهٔ KL در
رابطه~\eqref{cmp:eq:ch3-R-I-decomposition} نتیجه می‌دهد:
\begin{equation*}
\boxed{
\sum_{c=1}^{K}\frac{\omega_c}{|C_c|}
\sum_{i\in C_c}
D_{\mathrm{KL}}\bigl(\mathscr Y_i\|\overline{\mathcal V}\bigr)
=R_{\mathrm{KL}}+I(\mathcal V).
}
\end{equation*}
پس پراکندگی کلی نسبت به توزیع مشترک به دو
جزء درون‌طبقه‌ای و میان‌طبقه‌ای تجزیه می‌شود.
اگر $U$ شناسهٔ نمونهٔ مبدأ در آزمایش بالا باشد،
$C_0$ تابعی از $U$ است. طبق قاعدهٔ زنجیره‌ای
اطلاعات متقابل \cite{CoverThomas2006}:
\begin{equation}\label{cmp:eq:ch3-R-I-chain-rule}
\begin{aligned}
\mathcal I(U;C_\xi)
&=\mathcal I(C_0;C_\xi)+\mathcal I(U;C_\xi\mid C_0)\\
&=I(\mathcal V)+R_{\mathrm{KL}}.
\end{aligned}
\end{equation}
اثبات این تعبیر در بخش~\ref{app:geometry-R-I-channel}
از پیوست~\ref{app:model-geometry} آمده است.
بنابراین $R_{\mathrm{KL}}$ اطلاعات دربارهٔ شناسهٔ
عضو مبدأ پس از دانستن طبقهٔ آن و $I$ اطلاعات
دربارهٔ خود طبقه را اندازه می‌گیرد.

کانال درون‌طبقه‌ای با ورودی شناسهٔ عضو و خروجی
برچسب همسایگی همان نوع کانالی است که در بخش
~\ref{app:geometry-kl-information-radius} از پیوست
هندسه برای شعاع اطلاعاتی معرفی شد. در آنجا توزیع
ورودی روی اعضای منتخب بهینه می‌شود و ظرفیت کانال
به دست می‌آید؛ اما در $R_c$ توزیع ورودی روی همهٔ
اعضای طبقه یکنواخت است و در $I$ ورودی کانال
طبقهٔ مبدأ با وزن‌های تعیین‌شده است.

صفر شدن $R_{\mathrm{KL}}$ مانع یکسان شدن مراکز
نیست؛ همچنین مثبت بودن $I$ تمایز تک‌تک مراکز
رقیب را تضمین نمی‌کند. برای مثال اگر دو مرکز
یکسان و مرکز سوم متفاوت باشد، با وزن‌های مثبت
$I>0$ است، اما تمایزی میان دو مرکز نخست وجود
ندارد. از آنجا که تصمیم مدل به مقایسهٔ ناهمانندی
با مراکز رقیب وابسته است، افزون بر این سنجهٔ کلی،
باید جدایی دوتایی خود مراکز را مستقیماً بررسی کنیم.

% ادامه: تعریف عمومی ناهمانندی مراکز، همانند تعریف عمومی تراکم در ۳.۷.
